%=====================================================================
\section{Conclusion}
\label{sec:conclusion}
This study tested seven hypotheses from the temporal-network literature on a contact, a
messaging and a citation network within a common framework. The importance of temporal
information depends on the question asked. In the citation network, the final aggregate
describes degree heterogeneity and connectivity, whereas preferential attachment, ageing
and the growth law require historical degrees, publication ages or cumulative snapshots.
Within the selected AMiner corpus, sublinear degree preference and age-dependent citation
choice coexist with superlinear edge growth, but the sampled effective diameter does not
shrink. The WSB model describes citation histories compactly without an established
forecast advantage over linear extrapolation, so descriptive fit, temporal mechanism and
predictive value must be assessed separately.

In the contact and messaging networks, structural metrics depend strongly on the window
width, and time ordering restricts reachability over a finite horizon. Burstiness exceeds
the timestamp-shuffled baseline robustly in CM, whereas in PS the excess is small and depends
on the observation interval and the event definition; memory is weaker than under the
baseline in both networks. The persistence of CM ties relative to the baseline reverses
within both observation halves and at 14-day windows. Static communities conceal the midday
mixing in PS, while CM shows no comparable reorganisation. Conditional on major outbreaks,
epidemic peaks in PS occur later than under both reference models at the primary setting,
whereas CM does not support a general slowdown. These hypothesis-specific effects cannot be
combined into a single ranking of network types without further assumptions. Several
datasets per network type, broader corpus coverage, finer-resolution citation dates and
predictive tests under alternative censoring assumptions would strengthen the conclusions.
